% Q06_gold_answers.tex — GRADER CALIBRATION FIXTURE (correct, paraphrased).
% A drop-in replacement for paper/Q06_answers.tex (spec §4.7, §6.12): every answer is
% correct, but the wording, the precision, the order of the arguments and the card's
% field order are deliberately far from the reference, so the grader must recognise
% correctness by substance (the rubric's probes), not by matching text. B1 carries four
% decimals and cross-checks with the count form of the Wilson interval; B2 counts the
% all-failing pairs in words; B4 sums the other tail; B5 orders its fields differently.
%
% EXPECTED[B1][recomputation][4]: correct; every row of Probe A, filled with these numbers, reads ok, and the (c) line shows a width ratio of 0.507 at factor 4. No finding.
% EXPECTED[B2][recomputation][4]: correct; every row of Probe B reads ok. No finding.
% EXPECTED[B3][correctness][4]: correct; Probe C's rows read ok, and (b) names the unit, the dependence of retries and the effect on n. No finding.
% EXPECTED[B4][recomputation][4]: correct; every row of Probe D reads ok (the upper-tail sum equals the lower-tail sum by symmetry). No finding.
% EXPECTED[B5][card-fidelity][4]: every field agrees with the pinned range (lines 672–699); Oracle correctly says the range does not spell out the success predicate beyond "proves"; Variance computes no interval it attributes to the card. No finding on any field.
% Expected per-problem grades: B1 = 4, B2 = 4, B3 = 4, B4 = 4, B5 = 4 (3 is tolerated only with a declared-gap note and no flaw).
% A calibration run FAILS if any problem scores <= 2 or if any finding claims a wrong step,
% a wrong count or a contradiction with the range.
\documentclass[11pt]{article}
\input{preamble}

\title{Q06 --- The machine and the verifier loop \\ \large Written answers}
\author{learner}
\date{}

\begin{document}
\maketitle

% ---------------------------------------------------------------------------
\section*{Problem B1 (hand-calculation)}\label{ans:B1}
\begin{answer}
\textbf{(a)} With $z^2 = 1.96^2 = 3.8416$: \\
$\hat p = 17/20 = 0.85$. \\
$z^2/n = 0.19208$, so the denominator is $1.19208$. \\
centre: $(0.85 + 3.8416/40)/1.19208 = 0.94604/1.19208 = 0.7936$. \\
under the root: $0.85 \cdot 0.15/20 + 3.8416/(4 \cdot 400) = 0.006375 + 0.0024010
= 0.008776$; $\sqrt{0.008776} = 0.093680$. \\
half-width: $1.96 \cdot 0.093680/1.19208 = 0.18361/1.19208 = 0.1540$. \\
interval: $(0.6396,\ 0.9476)$ \quad width: $0.3080$. \\
Cross-check in counts: centre $= (17 + z^2/2)/(20 + z^2) = 18.9208/23.8416 = 0.7936$,
the same.

\textbf{(b)} Wald: $0.85 \pm 1.96 \cdot 0.079844 = 0.85 \pm 0.1565 = (0.6935,\ 1.0065)$. \\
The property it breaks: an interval for a proportion has to stay inside
$[0, 1]$, and this one ends at $1.0065$. \\
How its shape differs from Wilson's: Wald is centred on $\hat p$ itself; the
Wilson interval is lopsided around $\hat p$ (it reaches $0.21$ below and only
$0.10$ above), its centre moved toward one half.

\textbf{(c)} Factor: four, i.e.\ about 80 goals at the same rate \quad because the
width is essentially $2z\sqrt{\hat p(1-\hat p)/n}$, which scales with
$1/\sqrt{n}$; the correction terms in $z^2/n$ only shrink as $n$ grows.
\end{answer}

% ---------------------------------------------------------------------------
\section*{Problem B2 (hand-calculation)}\label{ans:B2}
\begin{answer}
\textbf{(a)} There are $\binom{5}{2} = 10$ ways to pick two of the five samples;
pass@2 is one minus the share of those pairs in which both samples fail. \\
t3 ($c = 4$): only one sample fails, so no pair fails: $\binom{1}{2} = 0$ and
pass@2 $= 1 - 0/10 = 1$. \\
t2 ($c = 2$): three samples fail, $\binom{3}{2} = 3$ failing pairs:
pass@2 $= 1 - 3/10 = 0.7$. \\
t1 ($c = 0$): every pair fails, $\binom{5}{2} = 10$: pass@2 $= 1 - 10/10 = 0$.

\textbf{(b)} pass@2 $= (1 + 0.7 + 0)/3 = 1.7/3 = 0.5667$ \quad pass@1 is each task's
share of successful samples, averaged: $(0.8 + 0.4 + 0)/3 = 0.4$

\textbf{(c)} Plug-in for t2: $1 - 0.6^2 = 0.64$ \quad above or below: below
($0.64 < 0.7$). \\
What the plug-in assumes: it treats $0.4$ as the true success rate and the two
attempts as fresh independent draws (so the same failing sample could be drawn
twice); the estimator only picks two different samples out of the five that
were actually run.
\end{answer}

% ---------------------------------------------------------------------------
\section*{Problem B3 (concept)}\label{ans:B3}
\begin{answer}
\textbf{(a)} Pass rate: one success among four first attempts, $0.25$ \quad
successes among all records: $3$ of $9$, $0.333$ \quad tasks solved at some
attempt: $3$ of $4$, $0.75$.

\textbf{(b)} What we want to estimate is a property of the prover on a
population of tasks, so the sample is the four tasks, not the nine rows. Rows
of one task are correlated (same task), and the policy creates them: a task
gets a second and third row only by failing, so the number of rows per task
depends on the outcome, and the 3/9 figure would change with the retry cap
alone. The 3/4 figure is solving within a budget of three tries, which is
another metric. Using all nine rows as $n$ pretends to have more independent
evidence than four tasks give, and the interval comes out too narrow.

\textbf{(c)} It must say how many attempts were allowed (3), whether a retry
used the earlier failure's feedback or was a fresh sample, how many tasks were
run, and what counted as solved; otherwise it matches neither a first-attempt
pass rate nor pass@3.
\end{answer}

% ---------------------------------------------------------------------------
\section*{Problem B4 (hand-calculation)}\label{ans:B4}
\begin{answer}
\textbf{(a)} Pass rates: A solves $14 + 1 = 15$, so $0.5$; B solves $14 + 10 = 24$,
so $0.8$ \\
The informative tasks: the 11 where exactly one prover succeeded; the
concordant 19 cannot favour either. \quad $b = 1$, $c = 10$

\textbf{(b)} With $X \sim \mathrm{Binomial}(11, 1/2)$, by symmetry I sum the
upper tail: $P(X \geq 10) = \bigl(\binom{11}{10} + \binom{11}{11}\bigr)/2^{11}
= (11 + 1)/2048 = 12/2048$, which equals $P(X \leq 1)$. Two-sided:
$p = 2 \cdot 12/2048 = 24/2048 = 0.0117$. \\
Decision at $\alpha = 0.05$: $0.0117 < 0.05$; the data reject equal
performance, in B's favour.

\textbf{(c)} In (a): the totals become 156 tasks, A $141/156 = 0.904$ and B
$150/156 = 0.962$, a gap of under six points. \quad In (b): unchanged,
$p = 0.0117$; the test sees only $b$ and $c$.
\end{answer}

% ---------------------------------------------------------------------------
\section*{Problem B5 (paper card)}\label{ans:B5}
\begin{answer}
\textbf{Source:} cheatsheet entry 02.4, lines 672--699 (the pinned range).

\textbf{Artifact:} the benchmark VeruSAGE-Bench (849 Verus proof tasks taken
from 8 verified Rust systems) and the engineered VeruSAGE agent; of the eight
systems only Atmosphere is safely held out.

\textbf{Task shape:} prove one function: the task file holds the function with
its spec unchanged, the types, spec functions and the signatures it depends on,
and no proof bodies, which the agent writes. Files are long (947 lines on
average, 496 of spec) with 2.4 helper lemmas and almost no loops.

\textbf{Metric:} percentage of the 849 tasks proved. Hands-off vs hands-on:
o4-mini 17 vs 41, GPT-5 50 vs 55, Sonnet 4 64 vs 58, Sonnet 4.5 81 vs 67;
AutoVerus 20. Side results: 83\,\% of Atmosphere, 16 of 17 tasks when starting
from partial human proofs (6 of 17 from scratch), and LLM proofs about 24 lines
at the median against 9 for people.

\textbf{Oracle:} Verus accepting the proof with the spec unchanged, plus a
cheat check by \texttt{verus-checker}; the range does not say more about the
success criterion than that a task is ``proved''.

\textbf{Cheat surface:} touching the spec, or proving nothing in a way the
verifier accepts; without the checker up to 14\,\% of results cheated, with it
under 1.5\,\%. The specs themselves are given and taken as correct.

\textbf{What is trusted:} the Verus toolchain and solver with \texttt{vstd};
the given specs and signatures (5 spec bugs turned up on Atmosphere, so this
trust is not free); the \texttt{verus log-all} extraction; the cheat checker.

\textbf{Budget:} hands-on: until verified, 20 minutes or 20 steps. Hands-off:
the range gives only Sonnet 4.5's cost, 7.2 minutes and USD 5.61 per task, and
no cap. The two settings differ in both harness and budget.

\textbf{Variance:} nothing: they seem to be single runs, so there is no
repeat-to-repeat spread, no seeds and no interval; and the hands-off/hands-on
gap mixes two changes at once.
\end{answer}

\end{document}
